% ---------------------------------------------------------------------
%  Chapitre 1 — Mesurer la dépendance qui menace le jugement
%  Principe du chapitre : décrire les données PUIS mesurer, AVANT tout
%  test, les trois canaux par lesquels les observations se ressemblent.
%  Réflexe du cours de sondages : connaître son plan d'échantillonnage
%  avant d'estimer.
% ---------------------------------------------------------------------
\chapter{Mesurer la dépendance qui menace le jugement}
\label{chap:donnees}

\section{La base des opérations}
\label{sec:donnees-trades}

Les méthodes présentées tout au long de ce rapport seront appliquées à une
base de données construite pour l'occasion. Celle-ci a été obtenue en
appliquant différentes stratégies à des données de marché passées. Ce procédé
porte le nom de test rétrospectif (\emph{backtesting}) et consiste à faire
rejouer une stratégie sur l'historique, comme si elle avait été employée à
l'époque, afin de vérifier si elle a su repérer un signal annonciateur de
l'évolution des cours.

Chaque application d'une stratégie donne lieu à une opération, c'est-à-dire à
l'achat d'un titre à une date donnée suivi de sa revente à une date
ultérieure. Cette opération, que la pratique désigne sous le nom de
\emph{trade}, constitue l'unité statistique de tout ce travail. Une ligne de
la base correspond donc à une position ouverte puis refermée par une
stratégie sur un titre déterminé.

Onze stratégies ont été retenues, dont dix stratégies techniques classiques
et un témoin. Ce témoin, que l'on nommera l'oracle, consulte les cours futurs
pour ne réaliser que des opérations gagnantes. Il ne s'agit évidemment pas
d'une stratégie utilisable, puisqu'elle exige de connaître l'avenir, mais
d'un instrument de contrôle. Son rôle est de vérifier que le protocole sait
reconnaître un avantage lorsqu'il existe réellement, et qu'il ne rejette pas
tout par excès de sévérité. Les dix autres stratégies, elles, décident
d'ouvrir une position à une date donnée en n'utilisant que les informations
disponibles avant cette date.

\begin{table}[H]
\centering
\caption{Les onze stratégies du banc d'essai et l'indicateur sur lequel chacune repose.}
\label{tab:strategies}
\renewcommand{\arraystretch}{1.25}
\small
\begin{tabular}{@{}ll@{}}
\toprule
Stratégie & Fondée sur \\
\midrule
\texttt{rsi\_classic}  & indicateur RSI (seuils 30/70) \\
\texttt{rsi\_strict}   & indicateur RSI (seuils 20/80) \\
\texttt{rsi\_trend}    & indicateur RSI avec filtre de tendance (MM200) \\
\texttt{ma\_crossover} & croisement de deux moyennes mobiles \\
\texttt{db\_bottom}    & figure en double creux \\
\texttt{dt\_top}       & figure en double sommet \\
\texttt{hs\_inverse}   & figure épaule-tête-épaule inversée \\
\texttt{hs\_classic}   & figure épaule-tête-épaule \\
\texttt{sr\_breakout}  & cassure d'une résistance \\
\texttt{sr\_breakdown} & cassure d'un support \\
\midrule
\texttt{oracle}        & témoin~: connaît le cours futur (triche) \\
\bottomrule
\end{tabular}
\end{table}

Ces stratégies reposent sur des indicateurs d'analyse technique dont il faut
dire un mot. L'indicateur RSI (\emph{Relative Strength Index}) compare
l'ampleur des hausses récentes d'un titre à celle de ses baisses, et produit
une valeur comprise entre 0 et 100. Une valeur basse signale un titre
fortement vendu, que les stratégies interprètent comme une occasion d'achat.
Une moyenne mobile, quant à elle, est la moyenne des cours des derniers jours,
recalculée chaque jour, qui lisse les soubresauts et laisse apparaître la
tendance de fond. Les figures évoquées dans le tableau, double creux ou
épaule-tête-épaule, sont des motifs graphiques que les praticiens croient
annonciateurs d'un retournement. Enfin, un support est un niveau de prix sous
lequel un titre a cessé plusieurs fois de descendre, une résistance le niveau
symétrique au-dessus duquel il a cessé de monter, et leur cassure est
interprétée comme le signe d'un mouvement durable.

J'ai choisi d'appliquer ces stratégies aux 501~titres composant le S\&P~500.
Cet indice regroupe les plus grandes sociétés cotées américaines et sert de
référence à l'ensemble du marché. Le fait qu'il compte 501 titres pour 500
sociétés tient à ce que quelques entreprises ont plusieurs catégories
d'actions cotées séparément. Un indice de ce type présente la particularité
d'être plus stable que chacune de ses composantes, puisque la chute d'un
titre isolé y est amortie par le comportement des autres, et sa tendance est
haussière sur la période retenue.

Pour qu'une stratégie soit validée, il faut qu'elle obtienne de meilleurs
résultats qu'une entrée au hasard, et que cet écart ne soit pas imputable au
seul échantillon dont nous disposons. Or la tendance est ici haussière, ce
qui nous ramène à la première difficulté annoncée en introduction. Comment
évaluer une stratégie lorsque le marché monte de lui-même~?

Pour répondre à cette question, on simule le hasard. À chaque opération
réalisée, on associe 200~opérations fictives de même exposition, portant sur
le même titre, pour la même durée et dans le même sens, mais dont la date
d'entrée est tirée au sort. La moyenne de leurs rendements définit le gain du
hasard pour cette opération.

On définit alors, pour chaque opération, son \edgevar{} comme le rendement
effectivement obtenu par la stratégie, éventuellement négatif, diminué de ce
gain du hasard,
\[
  \edgevar = \text{gain de la stratégie} - \text{gain du hasard}.
\]
Le mot anglais \edgevar{} désigne dans le vocabulaire du domaine l'avantage
qu'une méthode possède sur le hasard. Il est conservé ici comme nom de
variable, car sa traduction française se confondrait avec le terme général
d'avantage employé dans le fil du texte. Un \edgevar{} positif signifie que
la stratégie a fait mieux qu'une entrée au hasard de même exposition, et donc
que son choix du moment apporte une information que la seule dérive du marché
n'explique pas.

On obtient ainsi une base contenant, pour chaque stratégie, entre 978 et
8\,513 opérations réalisées entre janvier~2010 et juillet~2026, soit un total
de 45\,008 opérations. À chacune sont associées les informations décrites
dans le tableau~\ref{tab:dico-trades}, à savoir le titre concerné, la
stratégie employée, le secteur d'activité, les dates d'entrée et de sortie,
la durée de détention, le régime du marché et le contexte au moment de
l'achat, puis le rendement obtenu et l'\edgevar{} correspondant.

\begin{table}[H]
\centering
\caption{Description des variables d'une opération, illustrées sur un cas réel
de la base (achat du titre \texttt{A} par la stratégie \texttt{rsi\_classic}).}
\label{tab:dico-trades}
\renewcommand{\arraystretch}{1.3}
\small
\begin{tabular}{@{}>{\raggedright\arraybackslash}p{2.4cm}>{\raggedright\arraybackslash}p{2.6cm}>{\raggedright\arraybackslash}p{5.6cm}l@{}}
\toprule
Variable & Nature & Signification & Exemple \\
\midrule
\texttt{ticker}       & qualitative (501 mod.) & le code du titre traité (une même action revient dans plusieurs opérations) & \texttt{A} \\
\texttt{strategy}     & qualitative (11 mod.)  & la stratégie ayant produit l'opération & \texttt{rsi\_classic} \\
\texttt{sector}       & qualitative (11 mod.)  & le secteur d'activité du titre & Santé \\
\texttt{entry\_date}  & date                   & date d'achat & 2011-08-11 \\
\texttt{exit\_date}   & date                   & date de revente & 2012-01-20 \\
\texttt{holding\_days}& quantitative           & durée de détention, en jours & $162$ \\
\texttt{regime\_entry}& qualitative            & régime du marché à l'entrée & baissier \\
\texttt{vol\_entry}   & quantitative           & volatilité du titre à l'entrée & $0{,}043$ \\
\texttt{rsi\_entry}   & quantitative           & niveau de l'indicateur RSI à l'entrée & $30{,}1$ \\
\texttt{dist\_ma200}  & quantitative           & écart à la moyenne mobile sur 200 jours & $-0{,}210$ \\
\texttt{return\_net}  & quantitative           & rendement net réalisé par la stratégie & $+0{,}187$ \\
\texttt{rand\_return} & quantitative           & rendement d'une entrée au hasard équivalente & $+0{,}078$ \\
\texttt{edge}         & quantitative           & avantage = \texttt{return\_net} $-$ \texttt{rand\_return} & $+0{,}109$ \\
\texttt{win}          & binaire                & opération gagnante ($1$) ou perdante ($0$) & $1$ \\
\bottomrule
\end{tabular}
\end{table}

Deux de ces variables méritent une attention particulière, car elles portent
en germe toute la difficulté du chapitre. Le code du titre est une variable de
regroupement naturelle, puisque plusieurs opérations portent sur la même
action. Les dates d'entrée et de sortie, quant à elles, délimitent une période
d'exposition qui s'étend dans le temps, si bien que deux opérations peuvent se
chevaucher. Ce sont là les deux ingrédients de la dépendance mesurée en
section~\ref{sec:dep-mesuree}.

\section{Les rendements sectoriels}
\label{sec:donnees-secteurs}

L'étude de la dépendance demande une seconde source de données, distincte de
la précédente. Là où la base des opérations décrit les positions prises par
les stratégies, celle-ci décrit le marché lui-même.

Le S\&P~500 se répartit en onze secteurs d'activité, de la santé à l'énergie
en passant par la technologie. Pour chacun d'eux, on construit une série de
rendements journaliers, obtenue en moyennant chaque jour les rendements des
titres qui le composent. Le rendement d'un titre sur une journée est la
variation relative de son cours entre l'ouverture et la clôture, exprimée en
pourcentage.

Ces séries serviront à mettre en évidence le facteur marché en
section~\ref{sec:dep-canal2}, c'est-à-dire la part des mouvements qui est
commune à tous les secteurs et que l'on retrouvera au cœur du deuxième canal
de dépendance.

\section{La structure de dépendance, mesurée}
\label{sec:dep-mesuree}

Les tests employés au chapitre~\ref{chap:juger} ne sont valides que si les
observations sont indépendantes les unes des autres. Chacune doit apporter un
témoignage neuf, comme le ferait un juré qui n'aurait consulté aucun de ses
voisins. Or nos données ne satisfont pas cette condition, puisqu'elles se
ressemblent par le titre, par la date et par la durée.

Cette situation porte un nom en théorie des sondages. Notre base n'est pas un
tirage au hasard d'individus indépendants, mais un tirage par grappes, une
grappe désignant ici l'ensemble des opérations portant sur un même titre. Le
cours de sondage impose de décrire le plan d'échantillonnage avant d'estimer
quoi que ce soit, et c'est précisément la démarche suivie ici. Les trois
canaux par lesquels la dépendance s'installe sont mesurés maintenant, avant
que le moindre test ne soit appliqué.

% --- Canal 1 : les grappes par titre ---------------------------------
\subsection{Canal 1 --- la dépendance de titre}
\label{sec:dep-canal1}

Deux opérations portant sur un même titre sont liées par la gestion de la
société et par l'image qu'elle renvoie à l'extérieur. Elles subissent les
mêmes chocs propres à l'entreprise, qu'il s'agisse de la publication de
résultats, d'une décision de direction ou d'un litige, et cela indépendamment
du moment où elles sont ouvertes. On appellera cet effet la dépendance de
titre. Nos données constituent alors un tirage par grappes, où chaque grappe
est un titre.

Pour la mesurer, on définit l'estimateur suivant,
\[
  \hat\rho \;=\; \frac{\mathrm{MSB}-\mathrm{MSW}}
                      {\mathrm{MSB} + (n_0 - 1)\,\mathrm{MSW}},
\]
où MSB (\emph{Mean Square Between}), la variation entre les titres, mesure la
dispersion moyenne des \edgevar{} d'un titre à l'autre, tandis que MSW
(\emph{Mean Square Within}), la variation à l'intérieur d'un titre, mesure la
dispersion moyenne des \edgevar{} entre les opérations d'une même grappe.

Deux cas se présentent alors. Si les opérations d'un même titre ne se
ressemblent pas davantage que celles de titres différents, la variation est
identique à l'intérieur des grappes et entre elles, et l'on obtient
MSB~$\approx$~MSW. Si elles se ressemblent, ce qui est le propre d'une
dépendance intra-titre, la variation à l'intérieur d'une grappe s'effondre
tandis que la variation entre grappes subsiste, et l'on obtient
MSB~$>$~MSW.

On pourrait se contenter du simple quotient de MSB par MSW, qui vaudrait~1
sous indépendance. Mais un écart à~1, qu'il soit de $0{,}1$, de~3 ou de~100,
ne se rattache à aucune échelle, et l'on ne saurait dire s'il traduit une
dépendance forte ou faible. La formule retenue ci-dessus a l'avantage de
ramener la mesure entre~0 et~1, ce qui autorise une lecture directe du type
« 13,4~\% du résultat est imputable à une dépendance de grappe ».

Le terme $n_0$ correspond à la taille moyenne des grappes, c'est-à-dire au
nombre d'opérations par titre, corrigée de l'inégalité de ces tailles. Il
pondère la variation intra-titre au dénominateur, ce qui produit la
normalisation recherchée. Plus les grappes contiennent d'opérations, plus
$n_0$ est grand, et plus il faut de contraste entre les titres pour conclure
à une dépendance.

Le tableau~\ref{tab:dep-mesures} donne un $\hat\rho$ quasi nul pour la
plupart des stratégies, la valeur la plus élevée revenant à
\texttt{rsi\_strict} avec $0{,}134$.

La dépendance de titre est donc négligeable pour presque toutes les
stratégies. Un $\hat\rho$ faible ne garantit pourtant pas l'indépendance des
opérations. Il écarte ce premier canal, les deux suivants restant entiers.

% --- Canal 2 : les vagues simultanées --------------------------------
\subsection{Canal 2 --- la dépendance de marché}
\label{sec:dep-canal2}

La deuxième dépendance est celle du moment. Des opérations ouvertes à la même
période peuvent être liées parce qu'elles traversent les mêmes séances de
bourse. Sur un intervalle donné, les cours de différentes actions évoluent
conjointement, et les opérations qui les concernent se trouvent alors liées
par le marché.

Cette dépendance ne s'observe pas directement sur les opérations, mais sur le
marché lui-même. Si les secteurs évoluent de concert, alors des opérations
ouvertes au même moment, quel que soit leur secteur, subissent la même
poussée. On utilise pour la mesurer les rendements sectoriels quotidiens
introduits en section~\ref{sec:donnees-secteurs}.

On calcule d'abord les corrélations entre secteurs pris deux à deux, puis
leur moyenne, et l'on obtient $0{,}70$. Ce chiffre élevé indique que les
secteurs sont largement corrélés. Lorsqu'un secteur bouge, le marché tend à
bouger dans le même sens, et réciproquement, lorsque le marché s'oriente, la
plupart des secteurs le suivent.

On applique ensuite une analyse en composantes principales aux onze secteurs,
de façon à faire ressortir la structure de ces corrélations
(figure~\ref{fig:acp-marche}). Cette méthode construit des axes qui résument
au mieux la variance des données, le premier captant la plus grande part
possible, le deuxième la plus grande part de ce qui reste, et ainsi de suite.

Le premier axe capte ici 73~\% de la variance totale, et tous les secteurs
progressent ensemble le long de cet axe. C'est le mouvement global du marché,
que la finance nomme le facteur marché. Le deuxième axe, nettement moins
important, sépare les secteurs cycliques, qui réagissent fortement à la
conjoncture, de ceux qui sont défensifs et dont l'activité varie peu quelle
que soit la situation. La technologie illustre le premier cas, puisqu'elle
dépend du pouvoir d'achat des ménages et des investissements des entreprises,
tandis que les services aux collectivités, électricité, eau et gaz,
illustrent le second, leur consommation demeurant stable en période faste
comme en période de crise.

\begin{figure}[H]
  \centering
  \includegraphics[width=0.78\linewidth]{etape3_acp.png}
  \caption{Analyse en composantes principales des onze rendements
  sectoriels~: variance expliquée par composante (gauche) et position des
  secteurs sur les deux premiers axes (droite).}
  \label{fig:acp-marche}
\end{figure}

Sur l'histogramme de gauche, la première barre écrase toutes les autres,
puisqu'elle capte à elle seule 73~\% de la variance contre moins de 8~\% pour
la deuxième. Sur le nuage de droite, tous les secteurs se situent du même
côté de l'axe horizontal, ce qui confirme qu'ils progressent ensemble, tandis
que l'axe vertical les répartit entre profils défensifs et cycliques.

Cette cause commune n'est donc pas une hypothèse mais une mesure, puisque
73~\% des mouvements sont portés par une force unique. Deux opérations
ouvertes au même moment subissent largement le même sort, ce qui les rend
dépendantes par le marché.

% --- Canal 3 : les chevauchements dans le temps ----------------------
\subsection{Canal 3 --- la dépendance temporelle}
\label{sec:dep-canal3}

La suite du rapport demandera de regrouper les opérations par mois. Or les
opérations durent, avec une moyenne de 123~jours et une médiane de 53~jours,
et plus de 78~\% d'entre elles dépassent un mois. Une même position, encore
ouverte, traverse ainsi plusieurs mois consécutifs. Deux mois voisins
partagent donc des opérations communes et se ressemblent d'autant. Même
regroupées par mois, les observations ne sont pas indépendantes, ce qui
conduit à mesurer la corrélation d'un mois au suivant.

On emploie pour cela l'autocorrélation d'ordre~1, notée $\hat\rho(1)$. Le
principe consiste à décaler la série mensuelle d'un mois, puis à corréler la
série avec elle-même ainsi décalée, ce qui mesure à quel point un mois
ressemble au précédent. Les valeurs obtenues figurent dans le
tableau~\ref{tab:dep-mesures}.

\begin{table}[H]
\centering
\caption{Les dépendances mesurées, stratégie par stratégie~:
corrélation intra-titre $\hat\rho$ (canal~1) et autocorrélation
mensuelle $\hat\rho(1)$ (canal~3).}
\label{tab:dep-mesures}
\small
\begin{tabular}{@{}lrr@{}}
\toprule
Stratégie & $\hat\rho$ intra-titre & $\hat\rho(1)$ mensuelle \\
\midrule
\texttt{db\_bottom}    & $0{,}000$ & $0{,}36$ \\
\texttt{dt\_top}       & $0{,}000$ & $0{,}38$ \\
\texttt{hs\_classic}   & $0{,}087$ & $0{,}15$ \\
\texttt{hs\_inverse}   & $0{,}001$ & $0{,}31$ \\
\texttt{ma\_crossover} & $0{,}003$ & $0{,}56$ \\
\texttt{rsi\_classic}  & $0{,}000$ & $0{,}23$ \\
\texttt{rsi\_strict}   & $0{,}134$ & $0{,}03$ \\
\texttt{rsi\_trend}    & $0{,}010$ & $0{,}20$ \\
\texttt{sr\_breakdown} & $0{,}051$ & $0{,}11$ \\
\texttt{sr\_breakout}  & $0{,}004$ & $0{,}05$ \\
\bottomrule
\end{tabular}
\end{table}

Le $\hat\rho(1)$ atteint $0{,}56$ pour \texttt{ma\_crossover}, sa valeur la
plus forte, alors qu'il tombe à $0{,}03$ pour \texttt{rsi\_strict}. La
dépendance temporelle est donc prononcée pour les stratégies dont les
positions durent longtemps, et faible pour celles dont les positions sont
brèves. Il faut se garder d'y voir une prévisibilité du marché. Ces valeurs
traduisent le chevauchement mécanique de positions qui se prolongent, non une
régularité exploitable.

% ---- Conclusion : bilan des trois canaux + transition ch. suivant ----
Les trois canaux de dépendance sont désormais mesurés, et aucun n'est
systématiquement élevé. La dépendance de titre reste souvent proche de zéro,
et l'autocorrélation temporelle s'efface pour les stratégies aux positions
courtes. Chacun se manifeste pourtant fortement pour certaines stratégies, la
dépendance de titre pour \texttt{rsi\_strict}, la dépendance de marché pour
toutes, la dépendance temporelle pour \texttt{ma\_crossover}. Le
chapitre~\ref{chap:juger} montrera qu'ils suffisent à renverser des verdicts.

Nous pouvons maintenant juger les stratégies en sachant exactement ce qui
menace les tests.
